%\section{Ottica Fisica}
\begin{frame}{Ottica Fisica}
    \begin{itemize}
        \item L'ottica geometrica fallisce quando $\lambda$ è \textbf{comparabile} alle dimensioni degli ostacoli
        \item  Emergono fenomeni che richiedono di trattare la luce come onda:
        \begin{itemize}
            \item Dispersione: scomposizione della luce bianca
            \item Diffrazione: la luce aggira gli ostacoli
            \item Interferenza: le onde si sommano o si annullano
            \item Polarizzazione: selezione della direzione di oscillazione
        \end{itemize}
    \end{itemize}
\end{frame}

%========================================================

\begin{frame}{Dispersione e Diffrazione}
    \textbf{Dispersione cromatica}
    \begin{itemize}
        \item $n$ dipende da $\lambda$: frequenze diverse vengono rifratte con angoli diversi
        \item Un prisma scompone la luce bianca nei colori dello spettro
        \item Fenomeno naturale: l'\textbf{arcobaleno}
    \end{itemize}
    \vspace{0.3cm}
    \textbf{Diffrazione}
    \begin{itemize}
        \item Un'onda che attraversa un'apertura di dimensioni $\sim \lambda$ si allarga a ventaglio
        \item I bordi delle ombre reali non sono mai perfettamente netti
        \item Impossibile per particelle classiche: prova diretta della natura ondulatoria
    \end{itemize}
\end{frame}

% Slide 23 - Interferenza e Polarizzazione
\begin{frame}{Interferenza e Polarizzazione}
    \textbf{Interferenza}
    \begin{itemize}
        \item Quando due onde coerenti si sovrappongono, le ampiezze si sommano
        \item \textbf{Costruttiva}: cresta + cresta $\rightarrow$ intensità massima
        \item \textbf{Distruttiva}: cresta + valle $\rightarrow$ intensità nulla (buio)
        \item Dimostrazione classica: esperimento della doppia fenditura (Young, 1801)
    \end{itemize}
    \vspace{0.3cm}
    \textbf{Polarizzazione}
    \begin{itemize}
        \item La luce naturale oscilla in tutte le direzioni perpendicolari alla propagazione
        \item Un \textbf{polarizzatore} seleziona una sola direzione di oscillazione
        \item Applicazioni: occhiali da sole polarizzati, display LCD, fotografia
    \end{itemize}
\end{frame}